\section{Conclusiones}
\label{sec:conclusiones}

\begin{table}[htbp]
    \centering
    \caption{Conclusiones del an\'{a}lisis de flexibilidad de la l\'{i}nea SIM-012.}
    \label{tab:conclusiones}
    \setcounter{itemcount}{0}
    \begin{tabularx}{\textwidth}{@{}NX@{}}
        \toprule
        \multicolumn{1}{c}{\textbf{Item}} & \textbf{Conclusion} \\
        \midrule
        & La tuber\'{i}a de succi\'{o}n del tanque TK Blowtank (SIM-012) cumple los criterios de esfuerzo de ASME B31.3-2016 en los nueve casos de carga evaluados: la evaluaci\'{o}n de cumplimiento de c\'{o}digo del modelo arroj\'{o} CODE COMPLIANCE EVALUATION PASSED. \\
        & El esfuerzo de c\'{o}digo m\'{a}ximo es \num{93808.6}~\si{kPa} en el nodo 130, caso 8 (EXP) L8=L3-L6, con una relaci\'{o}n de utilizaci\'{o}n de 34.8~\% frente al admisible de \num{269891.8}~\si{kPa}; la condici\'{o}n que gobierna el dise\~{n}o es el rango de expansi\'{o}n t\'{e}rmica. \\
        & Los desplazamientos m\'{a}ximos del sistema son 6.0~\si{mm} (componente horizontal $DX$, casos con $T_2$ y de expansi\'{o}n) y las cargas en soportes y anclas son moderadas (carga vertical m\'{a}xima \num{21559}~\si{N} en el soporte del nodo 90), por lo que la configuraci\'{o}n de soporter\'{i}a del isom\'{e}trico es adecuada. \\
        & La comparaci\'{o}n de las cargas en la boquilla del tanque con los admisibles del fabricante del equipo queda pendiente como verificaci\'{o}n externa, por no disponerse de dichos admisibles a la fecha. \\
        \bottomrule
    \end{tabularx}
\end{table}

En s\'{i}ntesis, la l\'{i}nea SIM-012 es t\'{e}cnicamente viable seg\'{u}n
ASME B31.3-2016 con un margen de utilizaci\'{o}n superior a dos tercios del
admisible. No se identifican riesgos de sobre-esfuerzo, interferencias por
expansi\'{o}n ni sobrecarga de soportes, por lo que no se requieren acciones
correctivas sobre la configuraci\'{o}n analizada.
